\documentclass[11pt]{article}
\usepackage{amsmath,amssymb,amsthm}
\newtheorem{theorem}{Theorem}
\title{Positive Residues from Antisymmetric Root Numerators}
\author{}\date{}
\begin{document}\maketitle
\paragraph{Status.}
This is a complete written one-node construction with exact rational checks.
It is not a universal nested decoder, a Lean theorem, or a head-complexity
improvement. The construction preserves the positive exterior geometry;
it does not change any definition in the original model.

\begin{theorem}[Positive root realization]
Let $q\geq1$, let
\[
 p_{0,0}<p_{0,1}<p_{1,0}<p_{1,1}<\cdots<p_{q-1,0}<p_{q-1,1}<T/2,
\]
and let $t_s\in[T-p_{s,1},T-p_{s,0}]$.
For every $\sigma\in\{-1,1\}^q$ there are strictly positive
$c_{s,0},c_{s,1}$ such that
\[
 \sigma_s\sum_{r=0}^{q-1}\sigma_r
 \left(\frac{c_{r,0}}{t_s-p_{r,0}}-
       \frac{c_{r,1}}{t_s-p_{r,1}}\right)>0
 \quad(0\leq s<q).
\]
No additive bias is needed.
\end{theorem}
\begin{proof}
For each $s<q-1$ where $\sigma_s\ne\sigma_{s+1}$ choose a point
$b_s\in(p_{s,1},p_{s+1,0})$, and let $P(t)$ be the monic product
of $t-b_s$ over these selected boundaries. Write $a=\deg P\leq q-1$.
There is a sign $\eta$ such that $\eta P$ has sign $\sigma_s$
throughout every closed pair interval $[p_{s,0},p_{s,1}]$.
Put
\[
 Q(t)=\prod_{r,e}(t-p_{r,e}),\qquad
 N(t)=\eta(t-T/2)P(t)P(T-t).
\]
The numerator has degree $2a+1\leq2q-1$, and satisfies
$N(T-t)=-N(t)$. At every pole in pair $s$, its sign is
$-\sigma_s$, because $P(T-p_{s,e})>0$.
The derivative $Q'(p_{s,0})$ is negative and $Q'(p_{s,1})$ is
positive. Define
\[
 c_{s,0}=\frac{N(p_{s,0})}{\sigma_s Q'(p_{s,0})},\qquad
 c_{s,1}=\frac{N(p_{s,1})}{-\sigma_s Q'(p_{s,1})}.
\]
Both numbers are strictly positive. Partial fractions give
\[
 \frac{N(t)}{Q(t)}=
 \sum_s\sigma_s\left(\frac{c_{s,0}}{t-p_{s,0}}
                      -\frac{c_{s,1}}{t-p_{s,1}}\right).
\]
At $t_s$, the factors $t_s-T/2$ and $P(t_s)$ are positive,
while $\eta P(T-t_s)$ has sign $\sigma_s$. Also $Q(t_s)>0$.
This proves the required inequalities.
\end{proof}

\paragraph{A linear parameter family.}
Every polynomial $N$ of degree at most $2q-1$ with
$N(T-t)=-N(t)$ can be written uniquely as
\[
 N(t)=(t-T/2)R(t(T-t)),\qquad \deg R\leq q-1.
\]
To see this, translate $t$ by $T/2$: the odd polynomial has the form
$xS(x^2)$, and $t(T-t)=T^2/4-x^2$. The preceding constructive
numerator is in this family. For fixed poles and root routing, every
residue is linear in the $q$ coefficients of $R$.
If all these residues are positive, the root signs at
$t_s=T-p_{s,0}$ hold automatically by antisymmetry, without any
additional inequalities. All lower-node comparisons, when present,
remain homogeneous linear inequalities in the same coefficients after
eliminating their block biases. This supplies a smaller exact feasibility
problem; its feasibility for arbitrary trees remains unproved.

\paragraph{A limitation on common pools within this family.}
Suppose two root words $\sigma,\tau$ use the same strictly positive
residues $c$, and both have numerators in the displayed antisymmetric
family. At every pole, the ratio of their numerator values is
$\tau_s/\sigma_s\in\{-1,1\}$. The transformed poles
$p_{s,e}(T-p_{s,e})$ are all distinct, since $t(T-t)$ is strictly
increasing for $t<T/2$. Consequently $R_\tau^2-R_\sigma^2$ has
$2q$ distinct roots and degree at most $2q-2$. It is identically zero,
so $R_\tau=R_\sigma$ or $R_\tau=-R_\sigma$. Positivity excludes
zero numerator values at the poles. Thus the two root words agree
or differ by a global sign. This only limits common-pool arguments
inside this special numerator family, not adaptive residues in general.

\paragraph{Verification and missing step.}
The companion checker exhausts all root words through six pairs at
three reflected query choices per pair. It verifies strict residue
positivity, every query sign, the partial-fraction identity at extra
rational points, and antisymmetry. These checks complement the proof;
they do not establish a larger-tree theorem.
The missing step is to choose one positive residue vector satisfying
all nested input-block constraints at once. A one-node construction
has no preservation guarantee for any child block. No independent
choices of the polynomials $R$ can be substituted for that shared-vector
requirement.
\end{document}
